\chapter{Edge Computing and Metrics Collection}
\label{ch:edge_computing}

\section{Edge Computing Paradigm}
\label{sec:edge_paradigm}

\subsection{Definition and Characteristics}
\label{subsec:edge_definition}

Edge computing brings computation and data storage closer to the sources of data
generation, rather than relying exclusively on centralized cloud infrastructure
\cite{shi2016edge, satyanarayanan2017emergence, varghese2016challenges, mao2017survey}.
In the context of distributed and microservice-based systems, this means processing
telemetry, logs, and service-level signals near the host or service instance where
they originate. The defining characteristics of edge computing are:

\begin{itemize}
    \item \textbf{Low Latency}: Processing data near the point of origin reduces the
    delay associated with transmitting information to remote cloud infrastructure before
    analysis or action is taken.

    \item \textbf{Bandwidth Efficiency}: Local preprocessing, filtering, and inference
    reduce the volume of raw data that must be transmitted over wide-area network (WAN)
    links.

    \item \textbf{Data Locality}: Keeping data near its source supports privacy
    preservation, operational control, and regulatory compliance.

    \item \textbf{Resilience and Autonomy}: Edge systems can continue to perform useful
    computation even under degraded or intermittent connectivity to centralized services
    \cite{satyanarayanan2017emergence}.
\end{itemize}

\subsection{Edge vs.\ Cloud Computing}
\label{subsec:edge_vs_cloud}

Edge and cloud computing are generally regarded as complementary paradigms
\cite{shi2016edge, deng2016optimal}. Cloud infrastructure provides elastic compute
capacity and centralized orchestration, while the edge emphasizes locality, reduced
latency, and decentralized processing.

\begin{table}[H]
    \centering
    \begin{tabular}{|p{0.22\textwidth}|p{0.34\textwidth}|p{0.34\textwidth}|}
        \hline
        \textbf{Dimension} & \textbf{Edge Computing} & \textbf{Cloud Computing} \\
        \hline
        Processing Location & Near the data source & Remote centralized infrastructure \\
        \hline
        Latency & Lower for local decisions & Higher due to network transit \\
        \hline
        Bandwidth Usage & Reduced via local preprocessing & Higher when raw data is centralized \\
        \hline
        Data Privacy & Greater locality and control & Data aggregated centrally \\
        \hline
        Compute Capacity & Often resource-constrained & Highly elastic and scalable \\
        \hline
        Fault Scope & Failures often remain localized & Central failures affect larger portions \\
        \hline
    \end{tabular}
    \caption{Comparison of edge and cloud computing paradigms}
    \label{tab:edge_vs_cloud}
\end{table}

In practice, real-time inference runs at the edge while long-term storage, global
coordination, and large-scale model training are handled centrally. Edge computing is
also a natural fit for federated learning, where model adaptation occurs across nodes
without centralizing the underlying raw data \cite{mcmahan2017communication,
kairouz2021advances}.

\subsection{Edge Computing in Distributed Systems}
\label{subsec:edge_in_ds}

Microservice architectures generate large volumes of fine-grained telemetry across many
independently deployed components \cite{burns2016borg, newman2021building,
dragoni2017microservices}, making them natural candidates for edge-oriented monitoring.
Edge-based analysis can be realized through several architectural patterns:

\begin{itemize}
    \item host-level monitoring agents,
    \item service-local collectors,
    \item sidecar processes or containers,
    \item embedded telemetry modules.
\end{itemize}

The common principle across these patterns is that at least part of the collection,
transformation, or inference pipeline executes near the monitored component, enabling
local anomaly detection, streaming analysis, adaptive control, and privacy-preserving
monitoring.

\section{Metrics Collection}
\label{sec:metrics_collection}

\subsection{What Metrics to Collect}
\label{subsec:metrics_what}

The effectiveness of anomaly detection depends strongly on the observability signals
available to the detection model. In distributed systems, collected metrics typically
fall into four broad categories:

\begin{itemize}
    \item \textbf{Resource Utilization}: CPU utilization, memory consumption, disk
    activity, filesystem pressure, and process-level resource usage, reflecting
    infrastructure and runtime stress.

    \item \textbf{Network Performance}: Throughput, latency indicators, connection
    counts, packet-level characteristics, and retransmission behavior---particularly
    important in communication-heavy microservice systems.

    \item \textbf{Application Performance}: Request rate, error rate, response time,
    queueing delay, and service-level indicators (SLIs) reflecting the externally
    visible health and responsiveness of a service.

    \item \textbf{Internal Service Metrics}: Thread counts, garbage collection
    behavior, cache performance, queue depths, and connection pool utilization,
    exposing internal contention not visible at the infrastructure layer.
\end{itemize}

Anomaly detection benefits from combining infrastructure-level and application-level
signals, since anomalous behavior may manifest as isolated metric deviations or as
changes in the relationships between multiple correlated variables \cite{sigelman2010dapper,
beyer2016site, o11y2022}.

\subsection{How to Collect Metrics}
\label{subsec:metrics_how}

Metrics collection follows an instrumentation and acquisition pipeline where signals
are exposed, sampled, transported, and stored \cite{prometheus2016, turnbull2018prometheus}.
Common collection approaches include:

\begin{enumerate}
    \item \textbf{Application-Level Instrumentation}: Services expose internal
    counters, gauges, histograms, or summaries via language-specific instrumentation
    libraries.

    \item \textbf{Infrastructure and Kernel-Level Monitoring}: System metrics are
    collected using host-level exporters or low-level tracing mechanisms such as eBPF,
    which enables efficient kernel and syscall level observability \cite{gregg2019bpf}.

    \item \textbf{Pull-Based Collection}: Monitoring systems periodically scrape metric
    endpoints exposed by services or agents.

    \item \textbf{Push-Based Collection}: Monitored components actively emit metric
    updates to downstream collectors or aggregation systems.
\end{enumerate}

In anomaly detection pipelines, collected samples are organized into temporally ordered
windows before being consumed by statistical or machine learning models, since temporal
context is frequently necessary to distinguish normal variability from abnormal behavior
\cite{malhotra2015lstm, hundman2018detecting}.

\subsection{Metric Preprocessing}
\label{subsec:metrics_preprocessing}

Raw operational metrics are often noisy, heterogeneous, and expressed on different
numerical scales. Preprocessing is therefore required before the data can be used
effectively in detection models. Common steps include:

\begin{enumerate}
    \item \textbf{Normalization or Standardization}: Features are rescaled to
    comparable numerical ranges, which is critical for neural network models sensitive
    to scale differences across input dimensions.

    \item \textbf{Temporal Window Construction}: Sequential observations are segmented
    into fixed-length windows to form model inputs suitable for time-series analysis.

    \item \textbf{Missing Data Handling}: Incomplete time series arising from scrape
    failures or service restarts are addressed through interpolation, imputation,
    masking, or exclusion of incomplete samples.

    \item \textbf{Noise Reduction and Smoothing}: Short-term volatility may be
    attenuated through smoothing or aggregation to emphasize persistent behavioral
    patterns.

    \item \textbf{Feature Transformation}: Derived features such as rates of change,
    rolling statistics, or inter-metric ratios may be computed to expose temporal or
    relational structure not evident in raw signals.
\end{enumerate}

Care must be taken to ensure that preprocessing does not inadvertently remove the very
patterns intended for anomaly identification.

\section{Local Anomaly Detection at the Edge}
\label{sec:edge_detection}

\subsection{On-Device Inference}
\label{subsec:edge_detection_concept}

On-device inference executes anomaly detection models near the telemetry source rather
than within a fully centralized pipeline. This approach offers several advantages:

\begin{itemize}
    \item reduced response time for local anomaly identification,
    \item lower network utilization due to reduced raw-data transmission,
    \item improved privacy through local processing,
    \item increased operational autonomy under intermittent connectivity.
\end{itemize}

These properties are particularly relevant where centralized streaming analysis is
constrained by latency, bandwidth, reliability, or privacy considerations
\cite{shi2016edge, satyanarayanan2017emergence}. Edge inference is also compatible with
federated learning, where model training occurs across nodes without centralizing the
underlying raw data \cite{mcmahan2017communication, kairouz2021advances}.

\subsection{Statistical Baseline Methods}
\label{subsec:edge_statistical}

Before considering more expressive machine learning models, anomaly detection is often
approached using classical statistical methods, which are computationally lightweight
and serve as useful baselines \cite{basseville1993detection, chandola2009anomaly}:

\begin{itemize}
    \item \textbf{Moving Average}: Tracks local central tendency over a rolling window
    to identify deviations from recent behavior.

    \item \textbf{Rolling Standard Deviation}: Measures short-term variability and
    supports dynamic thresholding based on local dispersion.

    \item \textbf{Z-Score}: Quantifies how far an observation deviates from a reference
    mean relative to the estimated standard deviation.

    \item \textbf{Exponentially Weighted Moving Average (EWMA)}: Assigns greater weight
    to recent observations, commonly used in process monitoring and change detection
    \cite{basseville1993detection}.
\end{itemize}

While simple and interpretable, these methods are limited in their ability to model
nonlinear dependencies, long-range temporal context, or multivariate interactions among
correlated metrics.

\subsection{Machine Learning Methods at the Edge}
\label{subsec:edge_ml}

Machine learning methods extend statistical approaches by learning richer
representations of normal and abnormal behavior. Commonly studied model families
include:

\begin{itemize}
    \item \textbf{Tree-Based Methods}: Isolation Forest isolates anomalies through
    recursive partitioning of the feature space \cite{liu2008isolation}.

    \item \textbf{Kernel-Based Methods}: One-Class SVM estimates a decision boundary
    enclosing normal observations in a kernel-induced feature space
    \cite{scholkopf2001estimating}.

    \item \textbf{Neural Network Methods}: Deep learning models---feedforward
    autoencoders, recurrent networks, and sequence autoencoders---learn nonlinear and
    temporal patterns from multivariate telemetry \cite{chalapathy2019deep, pang2021deep}.
\end{itemize}

In sequential telemetry domains, temporal models are often advantageous because many
anomalies are defined not by isolated values, but by unusual progressions or deviations
from learned sequence dynamics.
